\clearpage
\thispagestyle{empty}
\pdfbookmark[0]{Beurtheilungen und Anzeigen nicht eigner Schriften}{beurtheilungen}
\SourceLeaf{486}
\begin{center}
\vspace*{.22\textheight}
{\Large BEURTHEILUNGEN UND ANZEIGEN\par}
\vspace{14mm}
NICHT EIGNER\par
\vspace{14mm}
{\Large S C H R I F T E N.\par}
\end{center}
\clearpage
\fancyhead[L]{\small ROHDE. UEBER DIE MASSEN DER PLANETEN.}
\SourceLeaf{488}
\begin{center}
{\large ROHDE. UEBER DIE MASSEN DER PLANETEN.\par}
\medskip
\rule{0.55\linewidth}{0.4pt}\\[2pt]
v. \textsc{Zach}. Monatliche Correspondenz. Band XII. Seite 413 . . 419. 1805. October.\\[2pt]
\rule{0.55\linewidth}{0.4pt}
\end{center}
\begin{center}Gutachten eines andern Recensenten über die Abhandlung des Hauptmann \textsc{Rohde} über die Massen der Planeten, aus einem Briefe desselben gezogen.\end{center}

\noindent --- --- --- Die --- --- --- Vertheidigung des Hauptmanns \textsc{Rohde}, die Sie mir gefälligst mitgetheilt haben, hat mir zwar als ein Muster einer, für eine schlechte Sache noch viel schlechtern Vertheidigung einigen Spass gemacht; indessen bedaure ich doch, dass ich gewissermassen mit Schuld bin, einen so undankbaren Streit veranlasst zu haben. Meine Hoffnung, dass eine Zurechtweisung den Hauptmann \textsc{Rohde} wenigstens vorsichtiger machen würde, ist nicht eingetroffen, und seine unüberlegte Antwort, die er so voreilig hat drucken lassen, ist wieder voll neuer Beweise seiner --- --- --- --- --- Nicht um Ihrer Antwort vorzugreifen (denn wahrscheinlich haben Sie diese bereits vollendet, und die Gründlichkeit und überall hervorleuchtende innige Selbstüberzeugung in jener Recension sind Bürge, dass niemand diese besser vertreten kann, als Sie selbst), sondern hauptsächlich weil Sie meine Meinung über diese Antikritik zu wissen verlangen, will ich mich noch einmal mit Ihnen über diese Sache unterhalten.

Meine Ansicht von diesem Streite ist folgende: Sie haben in Ihrer Recension, ohne bestreiten zu wollen, dass \textsc{Rohde}'s Methode, die Planeten-Massen zu bestimmen, in abstracto oder blos theoretisch genommen, richtig und zuverlässig sein könne, behauptet und bewiesen,

1) dass \textsc{Rohde} sie bei der Ceres und Pallas ganz falsch angewandt habe, und also schon um deswillen seine Resultate dafür gar nichts werth sein können;

2) dass es bei dem heutigen Zustande der practischen Astronomie nicht möglich ist, die mittlern Abstände der Planeten so genau zu bestimmen, als zu dieser Methode erforderlich wäre, ausser wenn man sie aus der Umlaufszeit mit Rücksicht auf ihre Massen berechnet.

Eine Vertheidigung, wie die des Hauptmanns \textsc{Rohde} gegen 1) ist mir doch noch nie vorgekommen; sie heisst mit andern Worten: \emph{`da meine Methode, wie der Recensent in 2) bewiesen hat, doch ganz unbrauchbar und ohne practischen Werth ist, so wäre es lächerlich gewesen, bei ihrer Anwendung noch besondere Sorgfalt aufzuwenden.'} Unglücklicher Weise hat er hierbei wieder einen neuen schülerhaften Fehler begangen; er schliesst so: das Vorrücken der Nachtgleichen beträgt in einem Jahre 0.01419 in Zeit, also in 1684 Tagen noch nicht 0.07 Tage. So etwas würde man einem Anfänger in der Astronomie nicht verzeihen: das Vorrücken der Nachtgleichen in einem Jahre ist 50$''$ und die Sonne braucht 0.014 Tage, um es nachzuholen, also braucht diese 0.07 Tage, um das Vorrücken von 1684 Tagen einzubrin\SourceLeaf{489}gen; aber die Ceres braucht dazu mehr, als 0.3 Tage, und so viel ist der siderische Umlauf länger als der tropische.

Alles, was \textsc{Rohde} gegen 2) vorbringt, ist durchaus nichts sagendes Geschwätz, was gar nicht zur Sache gehört; er meint, da die Beobachtungskunst heutiges Tages so sehr vervollkommnet sei, dass die sechste Decimale bei den mittlern Abständen als zuverlässig angesehen werden könne; allein es ist eine bekannte Sache, dass alle in den astronomischen Werken vorkommende Angaben von mittlern Distanzen der Planeten sowohl als des \textsc{Halley}'schen Cometen aus den Umlaufszeiten berechnet sind, entweder mit oder ohne Rücksicht auf die Massen. Die in \textsc{La Place} \emph{Exposition du Système du Monde} etc. befindlichen Angaben durfte also \textsc{Rohde} gar nicht brauchen, und es ist demnach eine leere Einbildung, wenn er behauptet, die Uebereinstimmung seiner Resultate bei denjenigen Planeten, die Trabanten haben, mit den hieraus gefundenen, bestätige die Gültigkeit seiner Methode.

Uebrigens scheint \textsc{Rohde} die Nichtigkeit seiner Antwort selbst gefühlt zu haben (das Gegentheil wäre auch, wenn man nicht die sonderbarsten Voraussetzungen annehmen wollte, unbegreiflich), denn er sucht einzulenken, und den Schein anzunehmen, als habe er gar nicht zur Absicht gehabt, bessere Bestimmungen der Planetenmassen zu liefern, sondern er habe blos eine gegen jeden Einwurf sichere, streng richtige Methode gesucht, ohne sich darum zu bekümmern, ob die Elemente und Data dazu jetzt schon mit hinreichender Genauigkeit durch die Beobachtungen gefunden werden könnten. Hierüber sind nun drei Bemerkungen zu machen:

\noindent I) \emph{Ist dieses Vorgeben nicht wahr}, denn der H. \textsc{Rohde} behauptet zu Ende seiner Abhandlung ganz deutlich, alle Perturbations-Rechnungen, die man bisher in Ansehung der Planeten Mercur, Mars und Venus angestellt hat, seien nichts werth, theils weil man ihre Massen ganz unrichtig angenommen, theils weil man die ohne Vergleich grossen Einwirkungen der Ceres, Pallas und Juno habe vernachlässigen müssen; ferner, er habe nicht durch langsame Mittheilung im Manuscript, sondern auf dem schnellen Wege des Druckes die Astronomen von seiner wichtigen Entdeckung benachrichtigen wollen, damit ja gleich alles, was rechnen könnte, sich vereinigte, mit seinen neuen Massen-Angaben die Perturbationsrechnungen ganz umzugestalten, \emph{pour accélérer de tous côtés la récolte, à laquelle les forces d'un seul homme ne suffisent pas, et les miennes point du tout} (gewiss das wahrste in der ganzen Abhandlung), \emph{vû la grande quantité de combinaisons analytiques nouvelles etc.}

\noindent II) Nicht blos, um den mittlern Abstand eines Planeten, ohne ihn aus der Umlaufszeit zu berechnen, sondern überhaupt, um irgend einen Abstand aus Beobachtungen zu finden, gibt es kein anderes Mittel, als dass man ihn in dem allbekannten Dreiecke, wo die drei Winkel die Commutation, Elongation und jährliche Parallaxe und dabei die Entfernung der Sonne von der Erde als bekannte Seite angenommen werden, trigonometrisch berechnet; die Winkel findet man durch die Länge der Sonne und durch die geocentrische und heliocentrische Länge des Planeten, letztere muss aus den Tafeln berechnet werden; bekanntlich kann man aber bei keinem Planeten die Richtigkeit dieser berechneten Länge auf 5$''$ und mehr verbürgen; eben so wenig ist die geocentrische beobachtete Länge der Sonne und der Planeten innerhalb ein paar Secunden erhalten; diese Angaben sind nur sehr geringe angeschlagen und leicht kann man hieraus für die verschiedenen Planeten die Resultate ziehen, die zeigen, dass die hieraus bei den Distanzen entstehenden Fehler gegen die erforderliche Genauigkeit enorm sind, und \textsc{Rohde}'s Methode ganz unbrauchbar machen.

Ist \textsc{Rohde} aus diesem Schlupfwinkel vertrieben, so wird er sich begnügen zu sagen, dass seine Methode doch theoretisch richtig sei, und dass doch vielleicht nach Jahrtausenden, wenn die Beobachtungen zehn oder hundertmal genauer sind, als jetzt, ihre Anwendung möglich sein würde. Dagegen bemerke ich nun

\SourceLeaf{490}
III) \textsc{Rohde}'s Methode ist auch theoretisch genommen, falsch, führt in unvermeidliche Zirkel und ist folglich ganz und gar unbrauchbar.

Ich habe schon vorhin bemerkt, dass, um den Abstand eines Planeten aus Beobachtungen zu finden, nothwendig der gleichzeitige Abstand der Sonne von der Erde erfordert wird, und je genauer jener sein soll, je genauer muss auch dieser bekannt sein; daher müssen nicht nur die Elemente des Sonnenlaufs aufs genaueste bekannt sein, sondern auch alle Störungen, die den Abstand der Sonne von der Erde afficiren, also namentlich auch die durch den Planeten selbst, dessen Abstand man bestimmen will, also (und dieses ist die \emph{Pointe} der Sache) muss dessen Masse schon bekannt sein; ist sie es nicht, so sind die Störungsgleichungen für den Abstand der Sonne von der Erde unvollständig, so ist dieser Abstand falsch, so ist der berechnete Abstand des Planeten von der Sonne falsch, so ist der mittlere Abstand falsch, der auf einen oder mehrere wahre gebaut ist. Ein Beispiel wird die Sache einleuchtend machen.

Gesetzt, es wollte jemand nach \textsc{Rohde}'s Methode die Masse des Mars bestimmen; um dessen mittlern Abstand aus Beobachtungen zu bestimmen, müsste er einen oder mehrere wahre haben, dazu braucht er die Sonnenabstände von der Erde. In unsern ältern Sonnentafeln fehlen nun aber, weil die Masse des Mars unbekannt, alle von diesem Planeten abhängige Gleichungen, (in Freiherrn von \textsc{Zach}'s neuesten Sonnentafeln Arg. IV. u. Arg. XV). Also blos wegen dieses Mangels würde der Abstand der Erde von der Sonne um $0.0000058 + 0.0000011 = 0.0000069$ fehlerhaft sein können, der daraus entspringende Fehler beim Abstand des Mars von der Sonne würde in demselben Verhältniss grösser sein, als dieser Abstand jenen übertrifft, also über $0.00001$ betragen. Man muss es dem Hauptmann \textsc{Rohde} selbst zu berechnen überlassen, welchen Einfluss ein solcher Fehler im Abstande auf die Masse haben würde; auch schweige ich davon, dass man, um nach \textsc{Rohde}'s Methode die Masse eines Planeten zu finden, nicht nur diese selbst schon haben müsste, sondern auch die aller andern Planeten, weil man sonst weder eine genaue Theorie der Erde haben könnte (um den Abstand der Sonne von der Erde zu berechnen) noch eine genaue Theorie des Planeten selbst, dessen heliocentrische Länge man aus den Tafeln nehmen müsste, und dessen wirkliche Abstände man nicht anders zur Findung des mittlern brauchen könnte, als wenn sie vorher von dem Einfluss aller fremden Perturbationen degagirt und so rein elliptisch gemacht wären. Doch genug und schon viel zu viel über dieses unreife Product.
